\subsection{Integral Domains and Fields of Quotients}

\begin{definition} \label{def:zero-divisor-integral-domain} \hlt{Zero Divisors and Integral Domains}\\
Let $R$ be a ring.
\begin{enumerate}[label=(\roman*)]
\setlength{\itemsep}{0pt}
\item A \hlt{zero divisor} is a nonzero $a\in R$ such that $\exists\,0\neq b\in R$ with $ab=0$ or $ba=0$.
\item An \hlt{integral domain} is a commutative ring with identity $1_R\neq0$ and no zero divisors.
\end{enumerate}
\end{definition}

\begin{proposition} \label{prop:integral-domain-cancellation} \hlt{Cancellation in Integral Domains}\\
Let $R$ be a commutative ring with identity $1_R\neq0$.\\
Then $R$ is an integral domain $\iff$ $(\forall r,x,y\in R)(r\neq0\land rx=ry\Rightarrow x=y)$.
\begin{proof}
\begin{enumerate}[label=\roman*.]
\setlength{\itemsep}{0pt}
\item[$(\Rightarrow)$] If $rx=ry$ and $r\neq0$, then $r(x-y)=0$; no zero divisors gives $x-y=0$.
\item[$(\Leftarrow)$] If $ab=0$ and $a\neq0$, then $ab=a0$; cancellation gives $b=0$, so no zero divisors exist.
\end{enumerate}
\hlt{Structural Commentary}: Domains are exactly cancellative commutative rings.\\
\hlt{Key Insight}: Zero divisors obstruct cancellation.
\end{proof}
\end{proposition}

\begin{definition} \label{def:field-of-quotients} \hlt{Field of Quotients}\\
Let $R$ be an integral domain. On $R\times(R\setminus\{0\})$, set $(a,b)\sim(c,d)$ $\Longleftrightarrow$ $ad=bc$.\\
The \hlt{field of quotients} of $R$ is $(R\times(R\setminus\{0\}))/\sim$; write $[(a,b)]=a/b$. For $b,d\neq0$, $\frac{a}{b}+\frac{c}{d}=\frac{ad+bc}{bd}$; $\frac{a}{b}\frac{c}{d}=\frac{ac}{bd}$.
\end{definition}

\begin{proposition} \label{prop:field-of-quotients-embeds-domain} \hlt{Domains Embed in Fields}\\
Let $R$ be an integral domain with field of quotients $K$.\\
Then $\iota:R\hookrightarrow K$; $\iota(a)=a/1$  is an injective ring homomorphism.
\begin{proof}
By Proposition \ref{prop:integral-domain-cancellation}, the relation and operations in Definition \ref{def:field-of-quotients} are well defined.\\
For $a,b\in R$, $\iota(a+b)=\frac{a+b}{1}=\frac{a}{1}+\frac{b}{1}$; $\iota(ab)=\frac{ab}{1}=\frac{a}{1}\frac{b}{1}$; and $\iota(1_R)=\frac{1_R}{1}$..\\
Thus $\iota$ is a ring homomorphism. If $\iota(a)=\iota(b)$, then $a/1=b/1$, so $a=b$; hence $\iota$ is injective.\\
\hlt{Structural Commentary}: The field of quotients formally adds denominators to a domain.\\
\hlt{Key Insight}: Integral domains are exactly the setting where fraction arithmetic is consistent.
\end{proof}
\end{proposition}
